\begin{ExerciseBlock}
\ExerciseHeading{05}{Spektrale Breite diskreter Linien}{quantum_optics_spectral_linewidth}
Eine Quecksilberdampf-Gasentladungslampe emittiert eine Linie bei $\lambda_0=\qty{546.1}{\nano\metre}$. Die Lebensdauer des oberen Zustands sei $\tau=\qty{2.5e-8}{\second}$, die Gas­temperatur $T=\qty{300}{\kelvin}$ und die Atommasse $m_\mathrm{Hg}=\qty{3.34e-25}{\kilogram}$. Der untere Zustand wird als langlebig angenommen. Vergleichen Sie die volle Breite bei halbem Maximum (FWHM) der natürlichen und der Dopplerverbreiterung.
\begin{enumerate}
\item Verwenden Sie $\Delta E_\mathrm{nat}=\hbar/\tau$. Bestimmen Sie $\Delta f_\mathrm{nat}$ in Hertz und die entsprechende Wellenlängenbreite $\Delta\lambda_\mathrm{nat}$.
\item Berechnen Sie die Standardabweichung $\sigma_{v_\parallel}=\sqrt{k_\mathrm B T/m_\mathrm{Hg}}$ der Geschwindigkeitskomponente in Beobachtungsrichtung. Bestimmen Sie mit der Dopplernäherung für $|v_\parallel|\ll c$ die volle Breite $\Delta\lambda_\mathrm{D}$.
\item Welcher Effekt dominiert bei Raumtemperatur?
\item Wie kann ein Laser ein schmaleres Spektrum erzeugen?
\end{enumerate}
\end{ExerciseBlock}
\begin{Solution}
\par\medskip\noindent\textbf{Musterlösung}\par\nopagebreak[4]
\begin{enumerate}
\item Für den exponentiellen Zerfall des oberen Zustands gilt $\Delta f_\mathrm{nat}=\Delta E_\mathrm{nat}/h=1/(2\pi\tau)\approx\qty{6.37}{\mega\hertz}$. Mit $\lambda=c/f$ folgt in linearer Näherung
\[
 \Delta\lambda_\mathrm{nat}=\frac{\lambda_0^2}{c}\,\Delta f_\mathrm{nat}\approx\qty{6.33e-6}{\nano\metre}.
\]
\item Die thermische Geschwindigkeitsstreuung beträgt $\sigma_{v_\parallel}\approx\qty{111}{\metre\per\second}$. Die lineare Dopplerverschiebung ergibt $\sigma_\lambda=(\lambda_0/c)\sigma_{v_\parallel}$. Für die Gaußkurve gilt
\[
 \Delta\lambda_\mathrm{D}=2\sqrt{2\ln 2}\,\sigma_\lambda\approx\qty{4.78e-4}{\nano\metre}.
\]
Die Dopplerbreite entspricht der in der \href{https://physics.nist.gov/Pubs/AtSpec/node20.html}{NIST-Übersicht zur Linienverbreiterung} angegebenen FWHM.
\item Das Verhältnis $\Delta\lambda_\mathrm{D}/\Delta\lambda_\mathrm{nat}\approx75$ zeigt, dass hier die Dopplerverbreiterung dominiert.
\item Der optische Resonator wählt Frequenzen in der Verstärkungsbande aus. Stimulierte Emission verstärkt das Feld in den ausgewählten Moden. Zusätzliche Frequenzselektion kann den Laser auf eine Mode begrenzen.
\end{enumerate}
\end{Solution}
